\documentclass[10pt,a4paper]{altacv}

\geometry{left=1.5cm,right=1.5cm,top=1.5cm,bottom=1.5cm}

\usepackage[utf8]{inputenc}
\usepackage[T1]{fontenc}
\usepackage[default]{lato}

\definecolor{VividPurple}{HTML}{78b3de}
\definecolor{SlateGrey}{HTML}{13161c}
\definecolor{LightGrey}{HTML}{6b5051}
\colorlet{heading}{VividPurple}
\colorlet{accent}{VividPurple}
\colorlet{emphasis}{SlateGrey}
\definecolor{DarkGrey}{HTML}{4D4D4D}
\colorlet{body}{DarkGrey}

\renewcommand{\itemmarker}{{\small\textbullet}}
\renewcommand{\ratingmarker}{\faCircle}

\addbibresource{sample.bib}

\begin{document}

\name{BEN HAZEM AMEL}
\tagline{Étudiante en Génie Logiciel et Systèmes d'Information}
\photo{4.5cm}{imagee.jpeg}
\personalinfo{
  \email{amelbenhazem@gmail.com}\\
  \phone{+216 51 019 482 / +216 29 391 923}\\
  \location{Ben Arous, Mornag}\\
  \linkedin{linkedin.com/in/amal-ben-hazem}\\
  \github{github.com/amelbenhazem}
}
\makecvheader

\cvsection{Profil}
Ingénieure logicielle en devenir, actuellement en dernière année du cycle d'ingénieur en Génie Logiciel et Systèmes d'Information à Tekup.
Passionnée par l'innovation et l'automatisation, je transforme les idées en solutions logicielles concrètes et performantes.
Curieuse et proactive, je maîtrise les technologies modernes du développement web, de la conteneurisation et du déploiement d'applications.
Je suis à la recherche d'un stage PFE stimulant où je pourrai mettre en pratique mes compétences en DevOps, CI/CD, Kubernetes et monitoring, tout en contribuant à des projets ambitieux et innovants.

\cvsection{Formation}

\cvevent{Cycle d'ingénieur en Génie Logiciel et Systèmes d'Information (cours du soir)}{Université Tekup}{Présent}{Ariana, Tunisie}
\cvevent{Licence en Technologies de l'Information et de la Communication (Télécommunications)}{ISTIC (Institut Supérieur des Technologies de l'Information et de la Communication)}{2023}{Borj Cédria, Tunisie}
\cvevent{Baccalauréat en Informatique}{Lycée Chebbi Mornag}{2018}{Mornag, Tunisie}

\cvsection{Expériences professionnelles}

\cvevent{Projet de fin d'année - AutoRemise – Application web de gestion bancaire}{Shamash Information Technology}{Juillet -- Août 2025}{Centre Urbain Nord, Tunis}
\begin{itemize}
  \item Développement d'une PWA automatisant la génération de PDF de remise bancaire à partir de fichiers Excel.
  \item Création de templates personnalisés pour plusieurs banques tunisiennes.
\end{itemize}
\textbf{Technologies :} Next.js, React, Tailwind CSS, Framer Motion, PDF-lib, XLSX, TypeScript, next-pwa

\vspace{1em}

\cvevent{Projet de fin d'études – Solution d'automatisation pour la préparation de la chimiothérapie}{ASQII}{Février -- Mai 2023}{Le Kram}
\begin{itemize}
  \item Développement d'une solution automatisée optimisant la préparation de traitements en chimiothérapie (UI Figma, Qt, C++).
\end{itemize}

\vspace{1em}

\cvevent{Projet de fin d'année – Armoire de stockage automatique sécurisée}{ASQII}{Juillet -- Septembre 2022}{Le Kram}
\begin{itemize}
  \item Développement d'une armoire de stockage automatisée et sécurisée (Qt, Arduino).
\end{itemize}

\vspace{1em}

\cvevent{Stage d'observation}{STEG}{Août -- Septembre 2020}{Tunis}
\begin{itemize}
  \item Observation des activités liées à la téléphonie, fibre optique et antennes radioélectriques.
\end{itemize}

\cvsection{Projets académiques}

\cvevent{}{Projet DevOps – Déploiement et monitoring d'une application Next.js}{2025}{Projet académique}
\begin{itemize}
  \item Déploiement Next.js conteneurisé sur k3s avec CI/CD GitHub Actions
  \item Gestion des secrets Supabase et reverse proxy Nginx
  \item Monitoring Prometheus/Grafana et dépannage Kubernetes (pods, rollout, scale)
  \item \textbf{Technologies : Docker, Kubernetes (k3s), GitHub Actions, Nginx, Supabase (PostgreSQL, Auth, Storage), Prometheus, Grafana, node-exporter, kube-state-metrics, DigitalOcean.}
\end{itemize}

\cvevent{}{Chat Application MERN avec CI/CD Docker/EC2}{2025}{Projet académique}
\begin{itemize}
  \item Développement d'une application de chat MERN (MongoDB, Express, React, Node.js) conteneurisée avec Docker.
  \item Orchestration multi-services (frontend, backend, base de données) via Docker Compose.
  \item Pipeline CI/CD avec GitLab pour build, tests automatisés et déploiement sur EC2.
  \item Gestion des secrets, tests unitaires frontend et débogage de containers.
  \item \textbf{Technologies : MERN Stack, Docker, Docker Compose, GitLab CI/CD, EC2, Yarn.}
\end{itemize}

\cvevent{}{Solution médicale web – e-Care}{2025}{Projet académique}
\begin{itemize}
  \item Plateforme web centralisant les données médicales, messagerie sécurisée et prise de rendez-vous.
  \item \textbf{Technologies : NestJS, Nuxt.js, PostgreSQL.}
\end{itemize}

\cvevent{}{Kanban Flow – Gestion de tâches}{2025}{Projet académique}
\begin{itemize}
  \item Tableau Kanban avec tâches en temps réel, suggestions IA et chatbot.
  \item \textbf{Technologies : React, Node.js, WebSocket, MongoDB.}
\end{itemize}

\cvevent{}{ImmoConnect – Plateforme immobilière}{2025}{Projet académique}
\begin{itemize}
  \item Gestion des annonces (vente/location), recherche avancée, planification des visites et dashboard pour agents.
  \item \textbf{Technologies : ASP.NET Core Web API, MySQL, Repository/Service Pattern.}
\end{itemize}

\cvevent{}{BookManager – Application de gestion de livres}{2024}{Projet académique}
\begin{itemize}
  \item CRUD, notation par étoiles, commentaires et tableau de bord administrateur.
  \item \textbf{Technologies : Angular, Guards, Reactive Forms.}
\end{itemize}

\cvevent{}{Restaurant Management System}{2024}{Projet académique}
\begin{itemize}
  \item Système complet de gestion de restaurant : menu, commandes, tables, personnel et clients.
  \item \textbf{Technologies : NestJS, GraphQL, TypeORM, MySQL, Swagger.}
\end{itemize}

\cvevent{}{Plateforme e-commerce agricole – Senyetna}{2024}{Projet académique}
\begin{itemize}
  \item Vente en ligne de produits agricoles avec avis clients et statistiques.
  \item \textbf{Technologies : Symfony, PostgreSQL, Chart.js.}
\end{itemize}

\cvsection{Compétences Techniques}

\textbf{Langages de Programmation :}
\cvtag{Java} \cvtag{Python} \cvtag{JavaScript} \cvtag{TypeScript}

\divider

\textbf{Frameworks \& Bibliothèques :}
\cvtag{Spring Boot} \cvtag{Symfony} \cvtag{Django} \cvtag{NestJS} \cvtag{Express.js}
\cvtag{React} \cvtag{Angular} \cvtag{Next.js} \cvtag{Nuxt.js} \cvtag{Ionic}
\cvtag{Bootstrap} \cvtag{Tailwind CSS} \cvtag{Swagger} \cvtag{Node.js} \cvtag{.NET} \cvtag{PHP}

\divider

\textbf{DevOps, Cloud \& Systèmes :}
\cvtag{Linux} \cvtag{Git} \cvtag{GitHub} \cvtag{GitLab CI/CD} \cvtag{Docker}
\cvtag{Kubernetes} \cvtag{Helm} \cvtag{Nginx} \cvtag{Prometheus} \cvtag{Grafana}
\cvtag{Terraform} \cvtag{Ansible} \cvtag{CI/CD Automation} \cvtag{Monitoring}
\cvtag{Cloud (AWS, DigitalOcean)} \cvtag{Virtualisation} \cvtag{Réseaux IP} \cvtag{Load Balancing}

\divider

\textbf{Tests \& Qualité Logicielle :}
\cvtag{JUnit} \cvtag{PyTest} \cvtag{Jest} \cvtag{Selenium} \cvtag{Automatisation des Tests} \cvtag{Qualité Logicielle} \cvtag{API REST}

\divider

\textbf{Big Data \& Intelligence Artificielle :}
\cvtag{Machine Learning} \cvtag{Deep Learning} \cvtag{Apprentissage Automatique}
\cvtag{Fouille de Données} \cvtag{Traitement d'Image} \cvtag{Recherche Intelligente} \cvtag{TensorFlow} \cvtag{PyTorch}

\divider

\textbf{Bases de Données :}
\cvtag{MySQL} \cvtag{PostgreSQL} \cvtag{SQL Server} \cvtag{SQLite} \cvtag{Oracle} \cvtag{MongoDB} \cvtag{Redis} \cvtag{Supabase}

\divider

\textbf{Architecture \& Design :}
\cvtag{MVC} \cvtag{MVP} \cvtag{MVVM}
\cvtag{Design Patterns} \cvtag{Clean Architecture}

\divider

\textbf{Web \& Marketing Digital :}
\cvtag{SEO} \cvtag{Analyse du Trafic Web} \cvtag{Optimisation de Contenu} \cvtag{Google Analytics} \cvtag{Social Media Analytics}

\cvsection{Langues}
\begin{itemize}
  \item \textbf{Arabe :} Langue maternelle
  \item \textbf{Français :} Courant
  \item \textbf{Anglais :} Courant
  \item \textbf{Allemand :} Niveau scolaire
\end{itemize}

\cvsection{Vie associative}

\cvevent{Club Croissant Rouge Tunisien ISTIC}{Volontaire et Responsable Communication et Marketing}{2020 -- 2023}{}
\cvevent{Club MakerLabs ISTIC}{Responsable Ressources Humaines}{2021 -- 2023}{}
\cvevent{ISTIC Google Club}{Membre}{2018 -- 2020}{}
\cvevent{Jeune Chambre Internationale Mornag}{Membre}{2020 -- 2021}{}
\cvevent{Spark ISTIC}{Coach Junior à la fondation BIAT pour la jeunesse}{2020 -- 2021}{}

\cvsection{Centres d'intérêt}
Sport \quad Musique \quad Lecture \quad Développement personnel \quad Cuisine

\end{document}
